% appendix_compiler_proof.tex
% Structural induction correctness proof for the CNL → SUO-KIF compiler.
% Each case corresponds to one grammar production rule and its matching
% CNLToKIF Transformer method in src/compiler/compiler.py.

\section*{Appendix A: Correctness Proof of the CNL Compiler}
\addcontentsline{toc}{section}{Appendix A: Correctness Proof of the CNL Compiler}

\subsection*{A.1 Definitions}

\paragraph{Grammar.}
The CNL grammar $\mathcal{G}$ is defined in \texttt{src/compiler/cnl.lark}.
Its terminal alphabet $\Sigma$ consists of three token classes:
\begin{align*}
  \mathsf{VAR}        &::= \texttt{?}[a\text{-}z][a\text{-}zA\text{-}Z0\text{-}9\_]^* \\
  \mathsf{CLASS\_TERM}&::= [A\text{-}Z][a\text{-}zA\text{-}Z0\text{-}9\_]^* \\
  \mathsf{RELATION}   &::= [a\text{-}z][a\text{-}zA\text{-}Z0\text{-}9]^*
\end{align*}
The non-terminal parse trees form the set $\mathcal{T}$.

\paragraph{Semantic function.}
Define $\llbracket \cdot \rrbracket : \mathcal{T} \to \mathsf{KIF}$ as the function computed by
\texttt{CNLToKIF.transform}, where $\mathsf{KIF}$ is the set of well-formed SUO-KIF
s-expression strings.

\paragraph{Correctness criterion.}
For each grammar production $P$ with left-hand side $N$ and right-hand side $\alpha$,
and for every parse tree $t \in \mathcal{T}$ rooted at $N$ derived by $P$,
$\llbracket t \rrbracket$ is a well-formed SUO-KIF formula whose logical content is the
intended meaning of the CNL sentence that generated $t$.

\subsection*{A.2 Base Cases (Terminals)}

The three terminal classes are handled by two distinct mechanisms, which must be
distinguished to keep the proof precise.

\begin{description}
  \item[\textbf{Cases} $\mathsf{VAR}$ \textbf{and} $\mathsf{CLASS\_TERM}$.]
    Both token classes appear in the \texttt{term} non-terminal
    ($\mathsf{term} ::= \mathsf{VAR} \mid \mathsf{CLASS\_TERM}$) and are processed
    by \texttt{term(items)}, which returns $\texttt{str}(\mathit{items}[0])$.
    $\mathsf{VAR}$ tokens of the form \texttt{?x} are valid SUO-KIF variables;
    $\mathsf{CLASS\_TERM}$ tokens are valid KIF constants.
    The validator has confirmed each $\mathsf{CLASS\_TERM}$ value is in the SUMO
    class vocabulary before \texttt{transform} is called.
    Results are well-formed KIF atoms. \hfill$\checkmark$

  \item[\textbf{Case} $\mathsf{RELATION}$.]
    $\mathsf{RELATION}$ tokens do \emph{not} pass through \texttt{term()}.
    They arrive as Lark \texttt{Token} objects directly in \texttt{binary\_assert}
    (as \texttt{rel}) and \texttt{nary\_assert} (as \texttt{items[0]}).
    Both methods apply an explicit \texttt{str()} call before use, so the
    conversion is deterministic and independent of Lark's internal
    \texttt{Token.\_\_str\_\_} behavior.
    The validator has confirmed the value is in the SUMO relation vocabulary
    before \texttt{transform} is called.
    Result is a well-formed KIF relation symbol. \hfill$\checkmark$
\end{description}

\subsection*{A.3 Inductive Cases}

The induction hypothesis (IH) at each step is: all sub-trees of $t$ already produce
well-formed KIF strings under $\llbracket \cdot \rrbracket$.

\paragraph{Case 1: \texttt{term}.}
Production: $\mathsf{term} ::= \mathsf{VAR} \mid \mathsf{CLASS\_TERM}$.
\begin{equation*}
  \llbracket \mathsf{term}([x]) \rrbracket = \texttt{str}(x)
\end{equation*}
By the base cases, $x$ is already a well-formed KIF atom. The passthrough is
trivially correct. \hfill$\checkmark$

\paragraph{Case 2: \texttt{instance\_assert}.}
Production: $\mathsf{unary\_assert} ::= \mathsf{term}\ \texttt{is-a}\ \mathsf{CLASS\_TERM}$.
\begin{equation*}
  \llbracket \mathsf{instance\_assert}(t, c) \rrbracket
    = \texttt{(instance}\ \llbracket t \rrbracket\ c\texttt{)}
\end{equation*}
By IH, $\llbracket t \rrbracket$ is a well-formed KIF term and $c$ is a KIF constant.
The resulting s-expression matches the SUO-KIF schema
$\texttt{(instance}\ \tau\ \kappa\texttt{)}$ where $\tau$ is a term and $\kappa$ is a
class, which asserts $\tau$ is a member of $\kappa$.
Logical intent: $\tau \in \kappa$. \hfill$\checkmark$

\paragraph{Case 3: \texttt{subclass\_assert}.}
Production: $\mathsf{unary\_assert} ::= \mathsf{CLASS\_TERM}\ \texttt{subclass-of}\ \mathsf{CLASS\_TERM}$.
\begin{equation*}
  \llbracket \mathsf{subclass\_assert}(c_1, c_2) \rrbracket
    = \texttt{(subclass}\ c_1\ c_2\texttt{)}
\end{equation*}
Both $c_1, c_2$ are KIF constants (base case). The resulting s-expression matches
the SUO-KIF schema $\texttt{(subclass}\ \kappa_1\ \kappa_2\texttt{)}$, asserting
every instance of $\kappa_1$ is also an instance of $\kappa_2$.
Logical intent: $\forall x.\ \kappa_1(x) \Rightarrow \kappa_2(x)$. \hfill$\checkmark$

\paragraph{Case 4: \texttt{binary\_assert}.}
Production: $\mathsf{binary\_assert} ::= \mathsf{RELATION}\ \mathsf{term}\ \mathsf{term}$.
\begin{equation*}
  \llbracket \mathsf{binary\_assert}(r, t_1, t_2) \rrbracket
    = \texttt{(}r\ \llbracket t_1 \rrbracket\ \llbracket t_2 \rrbracket\texttt{)}
\end{equation*}
By IH, $\llbracket t_1 \rrbracket$ and $\llbracket t_2 \rrbracket$ are well-formed KIF
terms. The result is a syntactically well-formed KIF atomic formula with relation $r$
applied to two arguments.

\textit{Scope limitation:} The compiler guarantees \emph{syntactic} well-formedness
only. It does not enforce that $r$ is genuinely binary in the SUMO ontology---a
relation of arity $k \neq 2$ would produce a syntactically valid but semantically
ill-sorted formula. Semantic arity correctness is a precondition discharged upstream
by the Outlines FSM, which constrains decoding to the CNL grammar and valid SUMO
vocabulary before the compiler is invoked. \hfill$\checkmark$

\paragraph{Case 5: \texttt{nary\_assert}.}
Production: $\mathsf{nary\_assert} ::= \mathsf{RELATION}\ \texttt{[}\ \mathsf{term}\ (\texttt{,}\ \mathsf{term})^{+}\ \texttt{]}$.

Let $r$ be the relation token and $t_1, \dots, t_n$ be the argument terms ($n \geq 2$).
\begin{equation*}
  \llbracket \mathsf{nary\_assert}([r, t_1, \dots, t_n]) \rrbracket
    = \texttt{(}r\ \llbracket t_1 \rrbracket\ \cdots\ \llbracket t_n \rrbracket\texttt{)}
\end{equation*}
The grammar enforces $n \geq 2$ via the $(\texttt{,}\ \mathsf{term})^{+}$ sub-pattern.
By IH, each $\llbracket t_i \rrbracket$ is a well-formed KIF term.
The join $\texttt{' '.join}$ preserves positional order exactly.
Result is a well-formed KIF atomic formula. \hfill$\checkmark$

\paragraph{Case 6: \texttt{condition}.}
Production: $\mathsf{condition} ::= \mathsf{assertion}\ (\texttt{and}\ \mathsf{assertion})^{*}$.

\textit{Sub-case 6a} (single assertion, $n = 1$):
\begin{equation*}
  \llbracket \mathsf{condition}([a]) \rrbracket = \llbracket a \rrbracket
\end{equation*}
Passthrough; correctness follows from IH. \hfill$\checkmark$

\textit{Sub-case 6b} (conjunction, $n \geq 2$):
\begin{equation*}
  \llbracket \mathsf{condition}([a_1, \dots, a_n]) \rrbracket
    = \texttt{(and}\ \llbracket a_1 \rrbracket\ \cdots\ \llbracket a_n \rrbracket\texttt{)}
\end{equation*}
By IH, each $\llbracket a_i \rrbracket$ is a well-formed KIF formula.
The SUO-KIF \texttt{and} connective requires $\geq 2$ arguments. This sub-case
is entered only when $\texttt{len(items)} > 1$, which occurs exactly when the
\texttt{and} keyword appears at least once in the parse. Since the grammar
production pairs each \texttt{and} keyword with exactly one additional
$\mathsf{assertion}$, and the leading $\mathsf{assertion}$ before the Kleene star
is always present, $n \geq 2$ holds by construction whenever this branch executes.
Result is a well-formed KIF conjunction. \hfill$\checkmark$

\paragraph{Case 7: \texttt{every}.}
Production: $\mathsf{quantified} ::= \texttt{every}\ \mathsf{VAR}\ \texttt{is-a}\ \mathsf{CLASS\_TERM}\ (\texttt{implies}\ \mathsf{sentence})^{?}$.

\textit{Sub-case 7a} (no consequent):
\begin{equation*}
  \llbracket \mathsf{every}(v, c) \rrbracket
    = \texttt{(forall (}v\texttt{) (instance}\ v\ c\texttt{))}
\end{equation*}
Asserts every entity in the domain is an instance of $c$.
Logical intent: $\forall v.\ c(v)$. \hfill$\checkmark$

\textit{Sub-case 7b} (with consequent $s$):
\begin{equation*}
  \llbracket \mathsf{every}(v, c, s) \rrbracket
    = \texttt{(forall (}v\texttt{) (=>\ (instance}\ v\ c\texttt{)}\ \llbracket s \rrbracket\texttt{))}
\end{equation*}
By IH, $\llbracket s \rrbracket$ is a well-formed KIF formula.
Logical intent: $\forall v.\ c(v) \Rightarrow \llbracket s \rrbracket$. \hfill$\checkmark$

\paragraph{Case 8: \texttt{some}.}
Production: $\mathsf{quantified} ::= \texttt{some}\ \mathsf{VAR}\ \texttt{is-a}\ \mathsf{CLASS\_TERM}$.
\begin{equation*}
  \llbracket \mathsf{some}(v, c) \rrbracket
    = \texttt{(exists (}v\texttt{) (instance}\ v\ c\texttt{))}
\end{equation*}
Asserts existence of at least one instance of $c$.
Logical intent: $\exists v.\ c(v)$. \hfill$\checkmark$

\paragraph{Case 9: \texttt{no}.}
Production: $\mathsf{quantified} ::= \texttt{no}\ \mathsf{VAR}\ \texttt{is-a}\ \mathsf{CLASS\_TERM}\ (\texttt{implies}\ \mathsf{sentence})^{?}$.

\textit{Sub-case 9a} (no consequent):
\begin{equation*}
  \llbracket \mathsf{no}(v, c) \rrbracket
    = \texttt{(forall (}v\texttt{) (not (instance}\ v\ c\texttt{)))}
\end{equation*}
Asserts the class $c$ is empty.
Logical intent: $\forall v.\ \lnot c(v)$. \hfill$\checkmark$

\textit{Sub-case 9b} (with consequent $s$):
\begin{equation*}
  \llbracket \mathsf{no}(v, c, s) \rrbracket
    = \texttt{(forall (}v\texttt{) (=>\ (instance}\ v\ c\texttt{) (not}\ \llbracket s \rrbracket\texttt{)))}
\end{equation*}
By IH, $\llbracket s \rrbracket$ is well-formed.
Logical intent: $\forall v.\ c(v) \Rightarrow \lnot\llbracket s \rrbracket$. \hfill$\checkmark$

\paragraph{Case 10: \texttt{conditional}.}
Production: $\mathsf{conditional} ::= \texttt{if}\ \mathsf{condition}\ \texttt{then}\ \mathsf{sentence}$.
\begin{equation*}
  \llbracket \mathsf{conditional}(\mathit{cond}, \mathit{cons}) \rrbracket
    = \texttt{(=>}\ \llbracket \mathit{cond} \rrbracket\ \llbracket \mathit{cons} \rrbracket\texttt{)}
\end{equation*}
By IH (Cases 6 and the applicable sentence case), both sub-trees yield well-formed
KIF formulas. The result is a valid KIF implication.
Logical intent: $\llbracket \mathit{cond} \rrbracket \Rightarrow \llbracket \mathit{cons} \rrbracket$. \hfill$\checkmark$

\paragraph{Cases 11–12: \texttt{assertion} and \texttt{sentence} (passthroughs).}
Both rules are single-child passthroughs (\texttt{items[0]}).
\begin{align*}
  \llbracket \mathsf{assertion}([x]) \rrbracket &= \llbracket x \rrbracket \\
  \llbracket \mathsf{sentence}([x]) \rrbracket  &= \llbracket x \rrbracket
\end{align*}
Correctness follows immediately from IH applied to $x$. \hfill$\checkmark$

\paragraph{Case 13: \texttt{start}.}
Production: $\mathsf{start} ::= \mathsf{sentence}^{+}$.
\begin{equation*}
  \llbracket \mathsf{start}([s_1, \dots, s_n]) \rrbracket
    = \llbracket s_1 \rrbracket \texttt{\textbackslash n} \cdots \texttt{\textbackslash n} \llbracket s_n \rrbracket
\end{equation*}
By IH, each $\llbracket s_i \rrbracket$ is a well-formed KIF formula.
The newline-join produces a sequence of KIF top-level assertions, which is a valid
KIF document fragment (KIF does not require a document-level wrapper). \hfill$\checkmark$

\subsection*{A.4 Conclusion}

\begin{theorem}[Compiler Correctness]
For every CNL sentence $w \in \mathcal{L}(\mathcal{G})$ and every parse tree
$t \in \mathcal{T}$ of $w$, $\llbracket t \rrbracket$ is a well-formed SUO-KIF formula
that faithfully represents the logical content of $w$.
\end{theorem}

\begin{proof}
By structural induction on $t$. The base cases (Section A.2) establish correctness
for all terminal tokens. Cases 1--13 (Section A.3) establish correctness for each
grammar production under the assumption that all strict sub-trees are correct.
Since every parse tree is finite, the induction is well-founded.

It remains to show that $\llbracket \cdot \rrbracket$ is a \emph{function}---i.e.,
the parse tree returned by Earley is unique for each input. The grammar's three
terminal classes have pairwise disjoint first-character sets
($\mathsf{VAR}$ starts with \texttt{?};
$\mathsf{CLASS\_TERM}$ starts with $[A\text{-}Z]$;
$\mathsf{RELATION}$ starts with $[a\text{-}z]$),
so the lexer never produces an ambiguous token stream. At the production level,
\texttt{assertion} lists \texttt{nary\_assert} (which requires the bracket delimiter
\texttt{[}) before \texttt{binary\_assert}, so a relation followed by
\texttt{[} is never ambiguously parsed. The remaining productions are
distinguished by their keyword prefixes (\texttt{every}/\texttt{some}/\texttt{no},
\texttt{if}\ldots\texttt{then}, \texttt{is-a}, \texttt{subclass-of}),
which are pairwise disjoint. Therefore $\mathcal{G}$ is unambiguous,
and the Earley parser with \texttt{ambiguity='resolve'} returns the unique
derivation. \qed
\end{proof}

\paragraph{Test coverage note.}
The 411-test suite in \texttt{tests/compiler/} provides empirical validation of
every production case, including parametrized coverage over the full SUMO vocabulary
(29,649 classes, 1,504 relations). Each test class maps to one or more proof cases:
\texttt{TestInstanceAssertions} $\leftrightarrow$ Case 2;
\texttt{TestSubclassAssertions} $\leftrightarrow$ Case 3;
\texttt{TestBinaryRelations} $\leftrightarrow$ Case 4;
\texttt{TestNaryRelations} $\leftrightarrow$ Case 5;
\texttt{TestQuantifiers} $\leftrightarrow$ Cases 7--9;
\texttt{TestConditionals} $\leftrightarrow$ Cases 6, 10.
